
Table~\ref{tab:data-partitions} summarizes the data partitions,
their use across experiments, and the corresponding appendix methods.

\begingroup
\fontfamily{phv}\selectfont
\fontsize{7.5}{8.5}\selectfont
\hypersetup{hidelinks}
\setlength{\tabcolsep}{3pt}
\renewcommand{\arraystretch}{1.08}

% Table 1 palette.
\definecolor{splitInk}{HTML}{26323D}
\definecolor{splitPrior}{HTML}{A18338}
\definecolor{splitFunctional}{HTML}{477FA5}
\definecolor{splitContext}{HTML}{438879}
\definecolor{splitInteraction}{HTML}{8266A2}
\definecolor{splitDesign}{HTML}{B56E58}
\color{splitInk}

% Uniform category shading; body text remains dark.
\newcommand{\splituse}[2]{%
  \cellcolor{split#1!8}\color{splitInk}#2}

% Link to individual rows using:
% \hyperlink{split:prior}{prior benchmark partition}
\newcommand{\splitname}[2]{%
  \hypertarget{split:#1}{\textbf{#2}}}

\begin{longtable}{@{}
  >{\raggedright\arraybackslash}
    p{\dimexpr.18\textwidth-6pt\relax}
  >{\raggedright\arraybackslash}
    p{\dimexpr.30\textwidth-6pt\relax}
  >{\raggedright\arraybackslash}
    p{\dimexpr.19\textwidth-6pt\relax}
  >{\raggedright\arraybackslash}
    p{\dimexpr.22\textwidth-6pt\relax}
  >{\raggedright\arraybackslash}p{.11\textwidth}
@{}}

\caption{\textbf{Data partitions and experiment locations.}
Counts denote sequence windows unless otherwise stated.
Background colors identify the experiment categories in
Table~\ref{tab:bend2-reference}; shading has no numerical meaning.
Appendix references distinguish experiment uses from
partition-construction methods.
Study-specific holdouts do not establish exclusion from
foundation-model pretraining.}
\label{tab:data-partitions}\\

\toprule
Partition / panel
& Fitting, selection and evaluation
& Main-paper uses
& Appendix uses
& Ref. \\
\midrule
\endfirsthead

\multicolumn{5}{l}{\tablename~\thetable\ (continued)}\\
\toprule
Partition / panel
& Fitting, selection and evaluation
& Main-paper uses
& Appendix uses
& Ref. \\
\midrule
\endhead

\bottomrule
\endfoot

\splitname{prior}{Prior benchmark}
& 1,024 fitting; 512 evaluation.
Markov tuning: 768/256 within fitting,
then refit on all 1,024.
& \splituse{Prior}{Prediction and calibration,
Fig.~\ref{fig:prior}}
& \splituse{Prior}{Calibration
(\ref{app:prior-repair});
mask placement and state exchange
(\ref{app:prior-routing}).
Donor analysis forms 256 pairs from the same 512 windows.}
& p.~\pageref{app:sampling_control}. \\

\addlinespace[3pt]
\splitname{shift}{Activation shift}
& 768 estimation; 128 selection;
1,536 confirmation windows from 48 blocks.
& \splituse{Prior}{Activation shift,
Section~\ref{sec:sequence-priors}}
& \splituse{Prior}{Shift, calibration,
direction controls and ensemble comparison
(\hyperref[app:prior-mode-decision]{\ref*{app:prior-routing}})}
& p.~\pageref{app:prior-mode-decision}. \\

\addlinespace[3pt]
\splitname{haplotypes}{Natural haplotypes}
& Evaluation only:
512 SNP sites, two carrier sequences per site.
Calibration uses the prior benchmark fitting set.
& \splituse{Prior}{Haplotype prediction,
Section~\ref{sec:sequence-priors}}
& \splituse{Prior}{Carried alleles and background sensitivity
(\ref{app:prior-memory});
calibration evaluation (\ref{app:prior-repair})}
& p.~\pageref{app:prior-memory}. \\

\addlinespace[3pt]
\splitname{opensplice}{OpenSplice folds}
& Five gene-held-out folds.
Label-efficiency analyses use nested training budgets.
& \splituse{Functional}{Mutation effects,
Fig.~\ref{fig:simulate}A}
& \splituse{Functional}{Model comparisons,
label efficiency and error checks
(\ref{app:simulate})}
& p.~\pageref{app:model-comparisons}. \\

\addlinespace[3pt]
\splitname{transfer}{Strict assay transfer}
& OpenSplice source fitting;
COMPASS evaluation on 88 genes.
Shared genes and overlapping sequences excluded.
& \splituse{Functional}{Cross-assay transfer,
Fig.~\ref{fig:simulate}B}
& \splituse{Functional}{Transfer evaluation
(\hyperref[app:model-comparisons]{\ref*{app:simulate}})}
& p.~\pageref{app:model-comparisons}. \\

\addlinespace[3pt]
\splitname{editing}{Editing-context folds}
& Five outer chromosome-group folds.
Tuning sensitivity uses three inner chromosome folds.
& \splituse{Functional}{Editing efficiency,
Section~\ref{sec:functional-prediction}}
& \splituse{Functional}{Editing-efficiency prediction
and tuning sensitivity
(\ref{app:editing-context})}
& p.~\pageref{app:editing-readouts}. \\

\addlinespace[3pt]
% Verify analysis-specific fold reuse before treating this
% group as a single shared assignment across all Siraj analyses.
\splitname{siraj}{Siraj readout folds}
& Element-grouped fitting and evaluation;
related allelic backgrounds kept together.
& \splituse{Context}{Cell contrasts,
Fig.~\ref{fig:context}}
& \splituse{Context}{Context effects
(\ref{app:context})}
& p.~\pageref{app:context},
p.~\pageref{app:interaction-ablation}. \\*
&
& \splituse{Interaction}{Background interactions,
Fig.~\ref{fig:compose}}
& \splituse{Interaction}{Interaction analysis and ablation
(\ref{app:interaction-ablation})}
& \\

\addlinespace[3pt]
\splitname{gtex}{GTEx genes}
& 12,146 training; 1,350 validation;
5,782 test genes.
Both analyses use the same test genes.
& \splituse{Context}{Tissue contrasts and identification,
Section~\ref{sec:biological-context}}
& \splituse{Context}{Both tissue analyses
(\ref{app:gtex-tissue-contrast})}
& p.~\pageref{app:gtex-tissue-contrast}. \\

\addlinespace[3pt]
\splitname{personal}{Personal-sequence folds}
& Two settings:
person-held-out;
person- and local-locus-held-out.
& Appendix only
& \splituse{Context}{Prediction at familiar and unseen loci
(\ref{app:personal-native})}
& p.~\pageref{app:personal-native}. \\

\addlinespace[3pt]
\splitname{compass-interaction}{COMPASS interactions}
& OpenSplice-trained readouts;
COMPASS interaction evaluation cohort.
Ablation reuses fixed predictions.
& \splituse{Interaction}{Interaction prediction,
Fig.~\ref{fig:compose}}
& \splituse{Interaction}{Prediction
(\ref{app:compass-predictions});
ablation (\ref{app:interaction-ablation})}
& p.~\pageref{app:compass-predictions}. \\

\addlinespace[3pt]
\splitname{adaptation}{Siraj adaptation}
& Locus-grouped training, validation and test;
68 training and 23 test loci.
& Appendix only
& \splituse{Interaction}{Conditional-response training
(\hyperref[app:conditional-development]{\ref*{app:interaction-ablation}})}
& p.~\pageref{app:conditional-development}. \\

\addlinespace[3pt]
\splitname{motifs}{Motif constructs}
& Train on singles/pairs;
evaluate triplets from held-out motif multisets.
& \splituse{Interaction}{Motif arrangements,
Section~\ref{sec:genetic-interactions}}
& \splituse{Interaction}{Triplet prediction
(\ref{app:triplets})}
& p.~\pageref{app:triplets}. \\

\addlinespace[3pt]
\splitname{response}{Siraj response matching}
& Response-matching evaluation cohort;
fitted readouts hold out complete elements.
& \splituse{Design}{Target matching,
Fig.~\ref{fig:control}}
& \splituse{Design}{Response-matching evaluation
(\ref{app:response-matching})}
& p.~\pageref{app:response-matching}. \\

\addlinespace[3pt]
\splitname{allele}{Siraj allele selectivity}
& Strength-control calibration excludes evaluation loci.
Repeat subsets provide measurement checks.
& \splituse{Design}{Allele-selective edits,
Section~\ref{sec:genomic-design}}
& \splituse{Design}{Selective mutation choice and controls
(\ref{app:final-siraj})}
& Cohort:
p.~\pageref{app:selective-candidates};
calibration:
p.~\pageref{app:selective-control}. \\

\addlinespace[3pt]
\splitname{cell}{Siraj cell selectivity}
& Five locus folds for native calibration;
fitted readouts supply held-out predictions.
& \splituse{Design}{Cell-selective edits,
Fig.~\ref{fig:control}C}
& \splituse{Design}{Cell-selective evaluation and controls
(\ref{app:cell-selective})}
& p.~\pageref{app:cell-selective-fitting}. \\

\addlinespace[3pt]
\splitname{atlas}{Atlas families}
& 4,223 training / 1,023 validation sequences;
refit on all 5,246.
Candidate families excluded from fitting.
& \splituse{Design}{Multi-cell selection,
Section~\ref{sec:genomic-design}}
& \splituse{Design}{Expression-program selection
(\ref{app:multicell-programming})}
& p.~\pageref{app:multicell-split}. \\

\end{longtable}
\endgroup